% !TEX root = ../main.tex
\section{Diseño de clases}\label{sec:clases}

El diagrama de clases de diseño traslada los conceptos del dominio a clases de software con atributos, operaciones, constructores y visibilidad. Es la base directa de la implementación: los nombres, las firmas y los tipos coinciden con los del código del repositorio. Cada clase del dominio indica como estereotipo el principio de asignación de responsabilidades que justifica su existencia o su papel principal \citar{larman2003uml} y, cuando participa en un patrón, el rol que cumple en él \citar{gamma2003patrones}.

\subsection{Convenciones del diagrama}

\begin{itemize}
  \item \textbf{Entidades y tipos de valor.} Las clases con identidad y estado mutable (partida, mesa, jugadores, mano, historial, estados, estrategias y jugadas) se relacionan mediante asociaciones. Las clases inmutables que se comparan por valor llevan el estereotipo «dataType», se implementan como \texttt{record} y se usan como tipo de atributo o de parámetro; por legibilidad se presentan en un diagrama aparte (Figura~\ref{fig:tipos-dato}).
  \item \textbf{Extremos de asociación.} El nombre de rol corresponde al campo o propiedad que implementa la asociación, con su visibilidad; la multiplicidad se indica en ambos extremos. Una multiplicidad \texttt{0..*} se implementa como colección y \{ordered\} indica que el orden es significativo.
  \item \textbf{Miembros.} Los miembros abstractos se muestran en cursiva y los estáticos subrayados. La visibilidad \texttt{\textasciitilde} (paquete) corresponde a \texttt{internal} de C\#: el miembro solo es accesible dentro de \texttt{Domino.Dominio}. Las subclases muestran los miembros que redefinen o implementan. Los tipos de valor posicionales no muestran su constructor, porque coincide con sus atributos.
  \item \textbf{Roles de patrón.} Los estereotipos \emph{Strategy}, \emph{ConcreteStrategy}, \emph{State}, \emph{ConcreteState}, \emph{Command}, \emph{ConcreteCommand}, \emph{Context}, \emph{Invoker} y \emph{Receiver} indican el papel de la clase en los patrones de la sección~\ref{sec:patrones}.
\end{itemize}

\diagrama*{diagrama_clases}{Diagrama de clases de diseño: entidades, interfaces y proyectos}{fig:diagrama-clases}

\diagrama*{tipos_dato}{Diagrama de clases de diseño: tipos de valor del dominio}{fig:tipos-dato}

\subsection{Clases del dominio}

{\tablacompacta
\begin{longtable}{@{}L{4.4cm}L{10.1cm}@{}}
  \caption{Clases del proyecto \texttt{Domino.Dominio} y su responsabilidad}\label{tab:clases}\\
  \toprule
  \cab{Clase} & \cab{Responsabilidad} \\ \midrule \endfirsthead
  \toprule
  \cab{Clase} & \cab{Responsabilidad} \\ \midrule \endhead
  \bottomrule \endlastfoot
  \texttt{ControladorJuego} & Recibe las operaciones del sistema que provienen de la interfaz, las delega en la partida, resuelve a continuación los turnos de las CPU y devuelve la vista del jugador humano. Genera la semilla de cada partida y coordina guardar y cargar. \\ \filasep
  \texttt{Partida} & Aplica las reglas del juego y coordina mesa, jugadores, reparto e historial. Delega en su estado actual las acciones que dependen de él (Context de State), crea las jugadas, las ejecuta, las registra y las revierte (Invoker de Command), y crea los jugadores con la estrategia de su nivel. \\ \filasep
  \texttt{Mesa} & Mantiene la cadena de fichas colocadas, valida si una ficha puede colocarse en un extremo, la orienta al colocarla y permite retirarla al deshacer. \\ \filasep
  \texttt{FichasParaRobar} & Entrega fichas al jugador que roba y las recibe de vuelta al deshacer, en orden inverso. \\ \filasep
  \texttt{HistorialJugadas} & Registra cada jugada en orden y extrae la secuencia que debe revertirse al deshacer. \\ \filasep
  \texttt{Repartidor} & Genera el juego de fichas del tipo de dominó, lo baraja con una semilla y lo reparte; entrega las fichas restantes como fichas para robar. \\ \filasep
  \texttt{Jugador} & Clase abstracta con el alias, la posición y la mano. Declara si el jugador es automático y si le aplica el límite de tiempo. \\ \filasep
  \texttt{JugadorHumano} & Jugador cuyas decisiones llegan desde la interfaz; le aplica el límite de tiempo. \\ \filasep
  \texttt{JugadorCpu} & Jugador automático: no le aplica el límite de tiempo y delega la decisión de jugada en su estrategia, que define su nivel. \\ \filasep
  \texttt{Mano} & Administra las fichas de un jugador y calcula su suma de puntos, su mayor doble y su ficha de mayor puntaje. \\ \filasep
  \texttt{IEstrategiaCpu} & Criterio con el que una CPU elige entre las opciones válidas. \\ \filasep
  \texttt{EstrategiaPrincipiante} & Elige al azar una de las opciones válidas. \\ \filasep
  \texttt{EstrategiaIntermedia} & Elige la opción cuya ficha suma más puntos, para descargar la mano; en empate prefiere los dobles. \\ \filasep
  \texttt{EstrategiaAvanzada} & Evalúa cada opción por el número de fichas propias que seguirían siendo jugables con los extremos resultantes y desempata por puntos. \\ \filasep
  \texttt{EstadoJuego} & Clase abstracta con el comportamiento que depende del estado: jugadas admitidas, opciones, control del tiempo y deshacer. Por omisión no admite ninguna acción. \\ \filasep
  \texttt{EstadoPreparada} & Solo admite iniciar: reparte y pasa a Salida. \\ \filasep
  \texttt{EstadoSalida} & Solo admite la ficha de salida, sin límite de tiempo, y pasa a EnCurso. \\ \filasep
  \texttt{EstadoEnCurso} & Admite colocar, robar y pasar según las reglas, controla el tiempo del turno y permite deshacer. \\ \filasep
  \texttt{EstadoTerminada} & No admite acciones; la vista muestra el resultado. \\ \filasep
  \texttt{Jugada} & Clase abstracta de las acciones del juego: valida, ejecuta y deshace la acción sobre la mesa, la mano y las fichas para robar, e indica si termina el turno. \\ \filasep
  \texttt{JugadaColocarFicha} & Quita la ficha de la mano y la coloca en el extremo indicado; al deshacer, la retira de la mesa y la devuelve a la mano. \\ \filasep
  \texttt{JugadaRobar} & Roba una ficha y recuerda cuál fue para devolverla al deshacer. No termina el turno. \\ \filasep
  \texttt{JugadaPasar} & Pasa el turno cuando no hay ficha para lanzar ni fichas para robar. \\ \filasep
  \texttt{JugadaTiempoAgotado} & Registra la pérdida del turno por tiempo. \\ \filasep
  \texttt{IReloj}, \texttt{IRepositorioPartidas} & Abstracciones propias para la hora actual y para guardar, cargar y listar partidas guardadas. \\
\end{longtable}}

{\tablacompacta
\begin{longtable}{@{}L{4.4cm}L{10.1cm}@{}}
  \caption{Tipos de valor del proyecto \texttt{Domino.Dominio}}\label{tab:tipos-dato}\\
  \toprule
  \cab{Tipo} & \cab{Responsabilidad} \\ \midrule \endfirsthead
  \toprule
  \cab{Tipo} & \cab{Responsabilidad} \\ \midrule \endhead
  \bottomrule \endlastfoot
  \texttt{Ficha} & Par de valores normalizado (\texttt{ValorA} $\leq$ \texttt{ValorB}), de modo que dos fichas iguales son iguales sin importar su orientación. Calcula sus puntos, si es doble y si contiene un valor. \\ \filasep
  \texttt{FichaColocada} & Ficha orientada en la mesa, con el valor que queda hacia cada lado. \\ \filasep
  \texttt{OpcionJugada} & Par ficha y extremo que representa una jugada válida disponible. \\ \filasep
  \texttt{PuntajeFinal} & Alias, puntos y fichas reveladas de un jugador al terminar la partida. \\ \filasep
  \texttt{Resultado} & Forma de cierre, puntajes de todos los jugadores y ganador; aplica el desempate de RN-09. \\ \filasep
  \texttt{ConfiguracionPartida} & Datos de CU-01 con su validación y los valores predeterminados; traduce el límite de tiempo a segundos. \\ \filasep
  \texttt{DefinicionTipoDomino} & Valor máximo, fichas por jugador y total de fichas de cada tipo de dominó. \\ \filasep
  \texttt{VistaPartida} & Información visible para un jugador: su mano, la cantidad de fichas de los demás, la mesa, las opciones válidas, las acciones permitidas, el tiempo restante, las últimas jugadas y, al terminar, el resultado. \\ \filasep
  \texttt{PartidaGuardada} & Configuración, semilla y jugadas de una partida: lo necesario para reconstruirla. \\
\end{longtable}}

En la versión 1.0, \texttt{Jugada} era un tipo de valor. Al aplicar Command pasó a ser una jerarquía de clases con comportamiento (sección~\ref{sec:patrones}). Sus operaciones de ejecución son de paquete, de modo que la interfaz, que recibe jugadas dentro de \texttt{VistaPartida}, puede leerlas pero no ejecutarlas.

\subsection{Comportamiento de las operaciones principales}

\begin{itemize}
  \item \textbf{Iniciar.} \texttt{Partida.Iniciar} delega en el estado. \texttt{EstadoPreparada} pide a la partida que reparta con el \texttt{Repartidor} y la semilla, determine quién abre y con qué ficha (RN-05) y pase a \texttt{EstadoSalida}. Si la salida corresponde a una CPU, \texttt{ControladorJuego} la juega de inmediato.
  \item \textbf{Jugadas.} \texttt{JugarFicha}, \texttt{Robar} y \texttt{Pasar} crean la jugada correspondiente y la entregan a \texttt{Procesar} del estado actual. \texttt{EstadoSalida} solo acepta la ficha de salida y pasa a \texttt{EstadoEnCurso}. \texttt{EstadoEnCurso} acepta la jugada si corresponde al jugador en turno y cumple sus propias condiciones (\texttt{EsValida}). En ambos casos, \texttt{Partida.Aplicar} ejecuta la jugada, la registra, verifica el fin y, si la jugada termina el turno, lo avanza. Robar no termina el turno: el jugador sigue robando hasta tener ficha para lanzar o agotar las fichas para robar (RN-07).
  \item \textbf{Turnos de las CPU.} \texttt{ResolverTurnosCpu} se ejecuta mientras el estado admita jugadas y el turno sea de una CPU: si hay opciones, la CPU decide con su estrategia sobre su propia vista; si no, roba o pasa.
  \item \textbf{Tiempo.} Cada cambio de turno registra \texttt{inicioTurno} con \texttt{IReloj}. Solo \texttt{EstadoEnCurso} controla el tiempo: si el jugador en turno tiene límite y se agotó, aplica una \texttt{JugadaTiempoAgotado} (RN-04). El controlador lo verifica antes de cada acción del humano, y la consola lo consulta periódicamente para que el turno se pierda en el momento en que se agota.
  \item \textbf{Fin.} \texttt{VerificarFin} declara el dominó cuando el jugador queda sin fichas (RN-08) y el juego cerrado cuando no quedan fichas para robar y ningún jugador tiene ficha para lanzar (RN-09); en ambos casos construye el \texttt{Resultado} y pasa a \texttt{EstadoTerminada}.
  \item \textbf{Deshacer.} Solo \texttt{EstadoEnCurso} lo admite. \texttt{HistorialJugadas.ExtraerHastaUltimaDe} entrega, de la más reciente a la más antigua, las jugadas desde la última del humano, y cada una se deshace sobre sus receptores. El turno vuelve al humano y, si la mesa quedó vacía, la partida regresa a \texttt{EstadoSalida} (RN-10).
  \item \textbf{Guardar y cargar.} \texttt{CrearGuardado} produce la configuración, la semilla y la lista de jugadas. \texttt{Restaurar} crea una partida con la misma configuración y semilla, la inicia, lo que reproduce el mismo reparto, y vuelve a procesar las jugadas guardadas en orden a través del estado, sin consultar a las CPU (RN-11).
\end{itemize}

\subsection{Clases fuera del dominio}

Las siguientes clases residen en los proyectos que dependen del dominio. Ninguna de ellas es referenciada por \texttt{Domino.Dominio}: se relacionan con él a través de \texttt{ControladorJuego} y de las interfaces que implementan. Cada proyecto ejecutable arma sus dependencias en su punto de entrada.

\begin{itemize}
  \item \texttt{Domino.Persistencia}: \texttt{RepositorioPartidasJson}, que guarda cada \texttt{PartidaGuardada} como un archivo JSON en una carpeta. Registra los cuatro tipos de jugada con un discriminador, de modo que el dominio no depende del formato del archivo.
  \item \texttt{Domino.Consola}: \texttt{InterfazConsola}, que lee las teclas y traduce cada acción en una operación del controlador; \texttt{PantallaConsola}, que dibuja el tablero completo en una sola pantalla con los puntos de cada ficha; y \texttt{RelojSistema}.
  \item \texttt{Domino.Simulacion}: \texttt{SimuladorPartidas} y \texttt{RelojSimulado}. El simulador ocupa el lugar del jugador humano, juega la primera opción válida, o roba o pasa cuando no la tiene, y en cada partida deshace, guarda y carga una vez.
\end{itemize}

\subsection{Tamaño del diseño}

El diseño suma 35 clases e interfaces en el dominio (26 entidades, interfaces y participantes de patrones, y 9 tipos de valor) y 6 en los demás proyectos, para un total de 41, sin contar las siete enumeraciones ni el proyecto de pruebas. La Tabla~\ref{tab:tamano} muestra de dónde provienen las 11 clases agregadas respecto de la Entrega~1.

\begin{table}[H]
\caption{Variación del número de clases respecto de la Entrega 1}\label{tab:tamano}
\tablacompacta
\begin{tabularx}{\textwidth}{@{}L{3.6cm}C{1.6cm}Y@{}}
  \toprule
  \cab{Cambio} & \cab{Clases} & \cab{Detalle} \\ \midrule
  Entrega 1 & 30 & Punto de partida. \\ \filasep
  Strategy & $+1$ & Se agregan \texttt{IEstrategiaCpu} y tres estrategias, y se retiran las tres subclases de \texttt{JugadorCpu}. \\ \filasep
  State & $+5$ & \texttt{EstadoJuego} y sus cuatro estados concretos. \\ \filasep
  Command & $+4$ & \texttt{Jugada} y sus cuatro jugadas concretas reemplazan el tipo de valor \texttt{Jugada}. \\ \filasep
  Consola & $+1$ & \texttt{PantallaConsola} separa el dibujo del tablero de la lectura de teclas. \\ \filasep
  Entrega 2 & 41 & \\
  \bottomrule
\end{tabularx}
\end{table}

El total supera el rango orientativo de 15 a 30 clases de la especificación. Cada clase adicional cumple un rol en un patrón justificado en la sección~\ref{sec:patrones}, y los participantes concretos son clases pequeñas: los estados y las jugadas concretas tienen entre una y cinco operaciones.
